% MoonLab poster palette sheet ("iiserb" color theme).
% Standalone document: builds with pdflatex OR lualatex, no extra fonts.
%   lualatex palette.tex   -> palette.pdf (2 pages)
% Page 1: swatches with names, HEX and RGB.
% Page 2: suggestions for using the palette cohesively.
\documentclass[a4paper,11pt]{article}
\usepackage[margin=2.2cm]{geometry}
\usepackage{xcolor}
\usepackage{tikz}
\usepackage{booktabs}
\usepackage{array}
\usepackage{hyperref}

% --- Theme colors (must match colorthemes/beamercolorthemeiiserb.sty) ---
\definecolor{spaceindigo}{HTML}{2B2D42}
\definecolor{glaucous}{HTML}{6A7FDB}
\definecolor{bubblegum}{HTML}{E76B74}
\definecolor{lemonchiffon}{HTML}{FBF2C0}
\colorlet{glaucouslight}{glaucous!12!white}
\colorlet{lemonlight}{lemonchiffon!35!white}
\colorlet{pinkdark}{bubblegum!75!spaceindigo}
\colorlet{glaucousdark}{glaucous!70!spaceindigo}
\definecolor{offwhite}{HTML}{FAFAFC}

\newcommand{\swatch}[1]{\tikz{\fill[#1] (0,0) rectangle (2.4,1.1); \draw[black!25] (0,0) rectangle (2.4,1.1);}}

\title{MoonLab Poster Palette \\ \large \texttt{iiserb} color theme}
\author{Multi-Robot Autonomy Lab, IISER Bhopal}
\date{}

\begin{document}
\maketitle

\section*{Core palette}
\begin{tabular}{@{}m{2.8cm}l>{\ttfamily}l>{\ttfamily}l@{}}
\toprule
 & \normalfont Name & \normalfont HEX & \normalfont RGB \\
\midrule
\swatch{spaceindigo}  & \normalfont Space Indigo  & 2B2D42 & 43, 45, 66 \\
\swatch{glaucous}     & \normalfont Glaucous      & 6A7FDB & 106, 127, 219 \\
\swatch{bubblegum}    & \normalfont Bubblegum Pink& E76B74 & 231, 107, 116 \\
\swatch{lemonchiffon} & \normalfont Lemon Chiffon & FBF2C0 & 251, 242, 192 \\
\bottomrule
\end{tabular}

\section*{Derived tints (defined in the theme file)}
\begin{tabular}{@{}m{2.8cm}ll@{}}
\toprule
 & \normalfont Name & \normalfont Mix \\
\midrule
\swatch{glaucouslight} & \normalfont glaucouslight & \texttt{glaucous!12!white} (example-block fill) \\
\swatch{lemonlight}    & \normalfont lemonlight    & \texttt{lemonchiffon!35!white} (alert-block fill) \\
\swatch{pinkdark}      & \normalfont pinkdark      & \texttt{bubblegum!75!spaceindigo} (pink text, if ever needed) \\
\swatch{glaucousdark}  & \normalfont glaucousdark  & \texttt{glaucous!70!spaceindigo} (links, small text) \\
\swatch{offwhite}      & \normalfont offwhite      & \texttt{FAFAFC} (reserved) \\
\bottomrule
\end{tabular}

\vfill
{\small Source of truth: \texttt{colorthemes/beamercolorthemeiiserb.sty}. If you change a color there, rebuild this sheet.}

\newpage
\section*{How to use the palette (keep it cohesive)}
Aim for roughly: \textbf{white 60\% / indigo 25\% / glaucous 10\% / pink + lemon under 5\%}.

\begin{description}
\item[Space Indigo — the workhorse.] Headline bar, block titles, headings,
bullets, plot axes, and the primary data series. The only dark color, so it
carries all text-like duties on white.
\item[White — the default.] Every block body. Do not tint body backgrounds;
print readability and contrast come first.
\item[Glaucous — the single recurring accent.] Headline rule, block
separators, example-block fills (12\% tint only), secondary plot elements.
\item[Bubblegum Pink — one highlight, used once.] The thin alert-block
separator and at most one contrasting plot series. Never a large fill, never
body text on white (use \texttt{pinkdark} if pink text is unavoidable).
\item[Lemon Chiffon — tint only.] Alert-block fill as a 35\% tint; never
lines or text on white — it is too light to read. The one exception is
footer links on the dark indigo bar, where full-strength lemon is readable.
\end{description}

\subsection*{Do}
\begin{itemize}
\item Pair one dark series with one accent series per figure
(\texttt{spaceindigo} vs \texttt{bubblegum}, or \texttt{spaceindigo} vs
\texttt{glaucousdark}); use \texttt{thick} / \texttt{very thick} lines —
posters are read from metres away.
\item Let the single pink element mark the single most important takeaway.
\item Keep tables black-on-white with \texttt{booktabs} rules.
\end{itemize}

\subsection*{Don't}
\begin{itemize}
\item Introduce extra colors (no pure \texttt{red}/\texttt{blue}/\texttt{green}
plot defaults — the template's placeholders already use theme colors).
\item Use more than 2--3 series per plot, or pink and lemon in the same figure.
\item Put lemon text on white, or white text anywhere except the indigo bars.
\end{itemize}
\end{document}
